%% kapa: exam=1 round=21 year=2010 period=1 form=A subject=civil_law subject_ko=민법 seq=11 number=11
%% source: Q-Net 감정평가사 기출문제 (cst00309) artlSeq=5024181 file=extracted/1008/5024181/2010년 제21회 감정평가사 1차시험 기출문제/1교시.hwp
%% conversion: kordoc 4.24.51 -> tools/gen1.py (auto); review=proofread-sonnet; flags=none
\begin{question}{11}
점유에 관한 설명으로 옳은 것은?
\begin{choices}
  \item 침탈자의 특별승계인이 악의인 경우, 점유자는 법정기간 내에 그 자에 대하여 점유물반환청구권을 행사할 수 있다.
  \item 점유자의 점유물반환청구권은 침탈사실을 안 날로부터 1년 내에 행사하여야 한다.
  \item 동산인 도품을 선의취득한 자에 대하여는 원소유자는 점유권에 기한 반환청구권을 행사할 수 있다.
  \item 선의의 점유자가 본권의 소에서 패소하면 그 판결이 확정된 때로부터 악의의 점유자로 본다.
  \item 선의의 점유자가 과실을 취득한 때에도 통상의 필요비의 상환을 청구할 수 있다.
\end{choices}
\end{question}
\answer{1}
